\bigskip\Needspace{12\baselineskip}
\pdfbookmark[1]{Gauß an Bode. 3. März 1803}{bode18030303}
\begin{center}
GAUSS AN BODE.\par
\medskip\hrule\medskip
Astronomisches Jahrbuch für 1806. S.\ 179—180. Berlin 1803.\par
\medskip\hrule
\end{center}\medskip
\begin{flushright}Braunschweig, 3. März 1803.\end{flushright}

Die wichtige Nachricht, dass der \textsc{Olbers}'sche Planet so früh in Lilienthal wiedergefunden wurde, werden Ew. — vermutlich schon von Herrn \textsc{Harding} selbst erhalten haben. Wahrscheinlich liegen Ihnen also auch die Beobachtungen von Dr.\ \textsc{Olbers} bereits vor. Ich hoffe, dass es Ihnen willkommen ist, auch Näheres über die Ergebnisse zu erfahren. Das Wetter ist hier seit dem 23. Febr. sehr schlecht. Jetzt kommt noch der Mondschein hinzu. Die Beobachtungen werden also wohl noch eine Weile unterbrochen bleiben.

Hier ist der Vergleich der Beobachtung von \textsc{Olbers} mit meinen Elementen V, nach denen die Ephemeride im Dezemberheft der M.\ C. berechnet wurde.
\begin{center}
\begingroup\small\setlength{\tabcolsep}{3pt}
\begin{tabular}{l|r|r|r|r|r}
& &\multicolumn{2}{c|}{Berechnete}&\multicolumn{2}{c}{Unterschied R—B}\\
1803&\shortstack{Mittl. Zeit\\in Bremen}&AR.&Declin.&AR.&Declin.\\\hline
Febr.\ 21&$17^{\mathrm u}\ 6^{\mathrm m}\ 10^{\mathrm s}$&$272^0\ 58'\ 47''$&$7^0\ 30'\ 40''$&$+2'\ 2''$&$-34''$\\
23&$15\ 24\ 36$&$273\ 31\ 14$&$7\ 45\ 4$&$+2\ 35$&$-57$
\end{tabular}\endgroup
\end{center}

Ich konnte dem Wunsch nicht widerstehen, auf dieser Grundlage eine vorläufige Verbesserung der Elemente vorzunehmen. Ich behalte mir vor, bessere Elemente anzugeben, sobald mehr und genauere Beobachtungen vorliegen. Ew. — werden aus dem Vergleich mit den Elementen V [S.\ 233 dieses Bandes] sehen, dass der Unterschied nur sehr gering ist:
\SourceLeaf{240}%
\begin{center}
\begingroup\small\setlength{\tabcolsep}{3pt}
\begin{tabular}{l r}
Epoche 1803 auf dem Seeberger Meridian $\ .\ .$&$221^0\ 28'\ 54''0$\\
Jährliche trop. Bewegung (365 T.) $\ .\ .$&$78\ 0\ 36.9$\\
Tägliche $\ .\ .\ .\ .\ .\ .\ .\ .\ .\ .\ .$&$769''416$\\
Sonnenferne 1803, gegenüber den Fixsternen ruhend $\ .$&$301^0\ 24'\ 13''$\\[3pt]
Knoten $\ .\ .\ .\ .\ .\ .\ .\ .\ .\ .\ .$&$172^0\ 28'\ 8''$\\
Exzentrizität $\ .\ .\ .\ .\ .\ .\ .\ .\ .$&$0.245619$\\
Logarithmus der halben großen Achse&$0.4426160$\\
Neigung der Bahn $\ .\ .\ .\ .\ .\ .\ .\ .$&$34^0\ 38'\ 20''$
\end{tabular}\endgroup
\end{center}

Diese Elemente geben die Beobachtung folgendermaßen wieder:
\begin{center}
\begingroup\small\setlength{\tabcolsep}{4pt}
\begin{tabular}{l|r|r|r|r}
&\multicolumn{2}{c|}{Berechnete}&\multicolumn{2}{c}{Unterschied R—B}\\
1803&AR.&Declin.&AR.&Declin.\\\hline
Febr.\ 21&$272^0\ 56'\ 25''0$&$7^0\ 31'\ 29''8$&$-20''0$&$+15''8$\\
23&$273\ 28\ 46.8$&$7\ 45\ 52.3$&$+7.8$&$-7.7$
\end{tabular}\endgroup
\end{center}
Man wird den Planeten also in diesem Jahr jederzeit so leicht finden können, wie man es sich nur wünschen kann.

Hoffentlich wird nun auch $\CeresSymbol$ bald wiedergefunden. Da sie diesmal tief am südlichen Himmel steht, wird die Wiederauffindung in Italien und Frankreich wohl früher gelingen als bei uns.

